\documentclass[12pt]{ctexart}
\author{Yan}
\usepackage{amsmath,amssymb}
\usepackage{geometry}
\usepackage{setspace}

\geometry{a4paper,margin=2.5cm}
\setstretch{1.25}

\title{炮棋规则书}
\date{2026.3.17}

\begin{document}

\maketitle
\tableofcontents
\newpage

%------------------------------------------------
\section{总则}

\subsection{游戏性质}

《炮棋》是一种在 $9\times9$ 方格棋盘上进行的双人原创策略棋类游戏。双方通过落子、升级、构成炮管、发动打炮、围捕吃子等机制争夺棋盘控制权，并在游戏结束时按棋子数量决定胜负。

\subsection{游戏核心}

炮棋的核心在于以下几个系统彼此交织、不断触发：
%不在于单纯地“放子”或“吃子”，而

等级系统：棋子有$1$级到$5$级之分，等级直接影响炮管形成、局部压制与吃子能力。

炮管系统：同等级连续棋子可以形成炮管发动打炮，对盘面造成冲击并升级。

围捕系统：玩家可以通过局部包围与总等级优势，实现对敌方棋子的吃子与转换。

连锁结算系统：一次落子后，可能连续触发己方与对方的打炮、吃子、再打炮、再吃子，直到双方都无法继续操作为止。

\subsection{游戏目标}

游戏结束时比较棋盘上双方棋子的数量，棋子数量较多的一方获胜。
蓝方由于后手，在最终结算时享有固定补偿。

%------------------------------------------------
\section{棋盘、坐标与棋子}

游戏由两名玩家进行，分别为红方与蓝方。双方各自操控己方全部棋子。

\subsection{棋盘与坐标}

游戏棋盘为 $9\times9$ 的方格棋盘，共 81 个格子。每个格子在任意时刻至多容纳一枚棋子，也可以为空格。

为统一描述棋盘位置，本规则采用平面坐标表示格子位置：
左上角为 $(1,1)$ ，右下角为 $(9,9)$ 。
即横坐标从左到右依次为$1,2,\dots,9$，纵坐标从上到下依次为$1,2,\dots,9$。

\subsection{棋子等级与上限}

棋子等级共有五级：1级、2级、3级、4级、5级。

棋子的等级最高为5级。若因打炮或其他规则效果使某枚棋子的等级超过 5，则该棋子强制记为 5级。这种封顶规则适用于所有由规则效果导致的等级增加，不论其原本是己方棋子、敌方棋子，还是由空格新生成的棋子。

5级棋子不能构成炮管，但仍然参与所有其他规则判定，例如被打炮影响、吃子判定、局部等级统计以及最终棋子数量统计。

%------------------------------------------------
\section{开局、回合与轮次}
游戏开始时，双方相对而坐。

\subsection{初始布局}

游戏开始时，按以下方式放置初始棋子：红方在 $(9,9)$ 放置一枚红方1级棋子，蓝方在 $(1,1)$ 放置一枚蓝方1级棋子。

换言之，从任意一方的视角看，自己右下角都是己方初始子，而自己左上角都有一枚对方的初始子。

\subsection{回合}

一个回合是指一名玩家从落子开始，到连锁结算全部完成的完整行动。

红方先进行回合，蓝方后进行回合，此后双方按顺序交替进行回合。

一个完整回合包含两个阶段：落子阶段、结算阶段。
具体为当前回合方先完成一次合法落子，然后进入由双方交替触发的结算过程。
只有当双方都无法再进行任何打炮或吃子操作时，该完整回合才结束。

结算过程遵循如下顺序：
\begin{itemize}
\item 先由当前回合方在 $\text{打炮阶段—吃子阶段}$ 之间反复循环，直到自己无法行动；
\item 然后轮到对方进行同样的 $\text{打炮阶段—吃子阶段}$ 循环，直到对方也无法行动；
\item 若一方的循环使得另一方获得了打炮或吃子的机会，则轮到另一方循环；
\item 如此往复，直到双方均无法继续操作。
\end{itemize}

双方均无法继续操作后回合结束。

\subsection{轮次}

一个轮次，是指红方完成一个完整回合，随后蓝方完成一个完整回合。整局游戏由若干个轮次构成。

%------------------------------------------------
\section{落子规则}

在落子阶段，玩家 必须且只能执行一次合法落子操作。
合法落子操作仅有以下两类：放置新棋子、升级已有棋子。

\subsection{放置棋子}

玩家可以在一个空格中放置一枚己方1级棋子，但该空格必须与至少一枚己方已有棋子边相邻。

这里的“边相邻”仅指上下左右四个方向，不包括四个对角方向。

\subsection{升级棋子}

玩家也可以选择对己方已有棋子进行升级，但仅允许以下两种情形：将一枚己方 1级棋子升级为 2级棋子，将一枚己方 2级棋子升级为 3级棋子。

以下升级在落子阶段是不允许的：将3级升到4级，将4级升到5级

因此，落子阶段所能达到的最高等级仅为3级，3级以上的棋子只能通过后续打炮等规则效果形成。

\subsection{无法落子的判定}

若轮到某一方进行回合时，该方既找不到任何与己方已有棋子边相邻的空格来放置己方1级棋子，也无法对任意己方 1级或 2级棋子进行合法升级，则该方视为 无法落子。

“无法落子”的判定只看该方在 落子阶段 是否存在合法操作，不考虑若进入结算阶段后本可通过打炮或吃子改变局面的可能性。

%------------------------------------------------
\section{炮管的定义与炮口判定}

\subsection{炮管形成}

均为某一等级k的不少于3枚的己方棋子在棋盘上连续相邻排列，则称其构成一门k级n长炮管，其中n为连续棋子数。

炮管等级最高为4级，因此5级棋子即使连续若干枚相连，也 不能构成炮管。

\subsection{最大长度原则}

对同一方向上的连续同级棋子段，只认定其 最大连续长度 所对应的炮管，不重复拆分出更短的炮管。
例如 $(2,3,3,3,3,1)$ 中只存在一门 3级4长炮管，而不额外再认定其中包含某门 3级3长炮管。

\subsection{炮口判定}

若某次盘面变化后新形成一门炮管，则考察构成这门炮管的所有棋子中，最后一个升级到该等级k的棋子。
若该棋子更靠近炮管的一端，则该端为炮口，炮口所朝向的方向即为该炮管的方向。
若它到炮管两端距离相同，则由当前执行该次操作的玩家任选一端作为炮口。

若新形成的某门炮管中，有多枚棋子是在同一次变化中同时达到该等级k的，则取这些棋子的坐标质心。
比较该质心与炮管两端点的距离，较近的一端为炮口；若两端距离相同，则仍由当前执行该次操作的玩家选择炮口方向。

若一次打炮后同时形成多门新炮，则每一门新炮的炮口方向，都应分别根据“在本次变化中达到该等级的棋子”独立判定，不同新炮之间互不干扰。

\subsection{炮管的表示方式}

一门炮管统一记作：$(k,n,A,\text{方},\text{向})$，
其中$k$为炮管等级，$n$为炮管长度，$A$为炮管坐标；
$\text{方}$为炮管方向，即横向还是纵向；
$\text{向}$为炮口方向，若为横向炮管则取左或右，若为纵向炮管则取上或下。

%------------------------------------------------
\section{打炮阶段与炮管集合}

\subsection{炮管集合}

某一方进入自己的打炮阶段时，应先扫描当前棋盘上该方所有符合条件的炮管，构成一个炮管集合。
炮管集合中应包含所有1级至4级的横向与竖向炮管。

\subsection{打炮流程}

在打炮阶段中，执行方按以下规则进行操作：

\begin{itemize}
\item 从当前炮管集合中任选一门炮，对其执行一次打炮后将其从炮管集合中移除；
\item 检查本次打炮后是否出现新的炮管，若出现新炮则按规则加入炮管集合；
\item 重复以上步骤，直到炮管集合为空。
\end{itemize}

在炮管集合中，炮的发射没有固定顺序，
只要某门炮已经在当前炮管集合中，执行方就可以任选一门发射，
不存在按加入顺序发射之类的限制。

若某门炮已经进入本次打炮阶段的炮管集合，则即使在后续盘面变化中这门炮原来的棋形已经不再完整，
或者其构成棋子等级发生变化，或者其所在结构不再符合当前盘面上的炮管条件，它仍然可以按照进入集合时所记录的五元组信息执行发射。
也就是说，已入集合的炮，不再重新验证其当前是否仍然是炮，它只有在被实际打出后，才会从集合中移除。

\subsection{新炮加入规则}

只要棋盘发生变化，就需要检查是否产生新的炮管。这既包括打炮造成的盘面变化，也包括吃子造成的盘面变化。

若一次打炮后同时形成多门新炮，则按以下顺序处理：
先识别本次打炮后出现的全部新炮，然后统一执行包含替代规则，最后将保留下来的新炮同时加入炮管集合。

所谓“包含替代规则”，是指若新形成的某门炮管A完全包含某门已记录炮管B的全部构成棋子，则保留新炮管 A，并将旧炮管B从炮管集合中移除。
该规则既适用于新炮与旧集合中已有炮之间的包含关系，也适用于同一次打炮所形成的多门新炮之间的包含关系。

%------------------------------------------------
\section{打炮的具体规则}

设当前执行方选择的炮管五元组为：$(k,n,A,\text{方},\text{向})$，其中$k$为炮管等级，$n$为炮管长度。

为描述打炮效果，对盘面格子临时采用如下数值记法：己方棋子按其等级记为正数、格记为0、敌方棋子按其等级记为负数。
这一正负记法仅用于描述打炮的变化规则。

\subsection{打炮效果}

打炮的效果分为前方攻击和内部升级两部分。

前方攻击效果为对于一门k级 n长炮管，从炮口出发，沿炮口方向向前数第1格到第n−2格，共n−2个格子，这些格子的数值均增加k。

内部升级效果为在该炮管本体内部，与炮口棋子的距离为奇数的所有棋子，等级增加1。

一次打炮中，炮口前方攻击与炮管本体奇数距离升级同时结算。
也就是说，应先根据打炮前记录的炮管信息确定所有受影响格子，再统一更新盘面，不得先结算其中一部分，再以中间结果影响另一部分。

若前方攻击效果要求增加棋盘外某些格子的数值，则这些越界部分无效。
但炮管本体内部奇数距离的升级，若对应棋子仍在棋盘内，则仍照常结算。

\subsection{效果落地规则}

打炮引起的“数值增加”按以下方式落地到盘面：

若某格原本为空格，记为0。打炮后若结果为正数，则在该格生成一枚己方对应等级的棋子；若结果仍为0，则该格保持为空。

若某格原本为敌方棋子，等级为m，记为−m。打炮后若结果仍为负数，则该格仍为敌方棋子，其等级为该负数的绝对值；
若结果恰为 0，则该格变为空格；若结果为正数，则该格改为己方棋子，等级为该正数。

若某格原本为己方棋子，等级为m，记为+m。打炮后若结果为更大的正数，则该格仍为己方棋子，等级按结果更新。

任何由打炮产生的等级增加，不论其作用于炮管本体上的己方棋子、炮口前方的空格、炮口前方的敌方棋子还是炮口前方的己方棋子，只要最终结果超过5，则该格最终强制记为5级。

一次打炮可能同时造成炮管本体内部某些棋子升级、前方某些空格生成新子、某些敌方棋子降级、消失或翻转为己方棋子，以及某些己方棋子升级。
上述变化都可能导致新炮形成，因此一次打炮后必须重新检查所有新炮。

%------------------------------------------------
\section{吃子规则}

炮棋中的吃子不是普通棋类中那种直接走到对方位置进行吃子，
而是通过周围己方棋子达到一定数量门槛且局部区域内己方总等级高于对方，
来实现对敌方棋子的围捕与转换。

\subsection{可吃棋子}

对于一枚敌方棋子，若其周围相邻格中的己方棋子数量不少于一半，则该敌方棋子成为可吃棋子。
这里周围相邻格采用八邻域，即包括上下左右和四个对角方向中所有实际存在的邻格。

根据敌方棋子所在位置不同，周围实际邻格数量不同，对应门槛如下：
若位于内部格，周围共有 8 个邻格，至少需要 4 枚己方棋子；
若位于边缘格但非角格，周围共有 5 个邻格，至少需要 3 枚己方棋子；
若位于角格，周围共有 3 个邻格，至少需要 2 枚己方棋子。

\subsection{等级比较}

若某枚敌方棋子已经满足数量门槛，则进一步比较其所在局部区域内双方棋子的等级总和。
若该棋子位于内部格，则取其所在九宫格；
若位于边缘格，则取其所在六宫格；
若位于角格，则取其所在四宫格。
其中六宫格、四宫格都指以该棋子所在位置为中心、实际落在棋盘内的局部区域。

在对应局部区域内，分别统计己方棋子的总等级数与对方棋子的总等级数。空格记为 0，且不属于任何一方。
若满足 $\text{己方总等级数}>\text{对方总等级数}$，则该枚敌方棋子可以被吃掉。

\subsection{吃子的结果}

若某枚敌方棋子被吃掉，则移除该敌方棋子，并在其原位置放置一枚己方 1级棋子。

若当前执行方进入吃子阶段后，同时存在多个可吃目标，则由执行方任选其中一个进行吃子。

每次进入吃子阶段时，执行方 至多吃掉一枚棋子。即使此时盘面上还存在其他可吃目标，也不能在同一次吃子阶段中继续连吃。
若执行方成功吃掉一枚棋子，则当前吃子阶段立即结束，棋盘既然已经发生变化，就必须重新进入该方的打炮阶段，重新检查是否出现新的炮管。

%------------------------------------------------
\section{单方结算与双方连锁结算}

\subsection{单方结算}

某一方进入自己的结算时，应按如下方式进行

\begin{itemize}
\item 结算自己的打炮阶段，然后进入吃子阶段；
\item 若吃子成功，则回到自己的打炮阶段；
\item 如此反复。
\end{itemize}

只有当某一方同时满足没有任何炮管可以继续打炮且没有任何敌方棋子可以继续吃掉时，该方的本轮结算才结束：

\subsection{双方结算的往复触发}

当当前执行方完成自己的整轮结算后，对方开始进行自己的结算，流程与之完全相同。

若对方的结算再次改变盘面，使当前回合方重新获得可打炮机会或可吃子机会，则当前回合方应再次开始自己的结算。

只要棋盘发生变化，就必须重新留意：是否产生了新的炮管、是否出现了新的可吃棋子、是否使原本不可继续的结算重新可以继续。

一次完整回合只有在双方都无法再进行任何打炮或吃子时才真正结束，轮到下一名玩家进入新的落子阶段。

%------------------------------------------------
\section{游戏结束与胜负判定}

游戏在以下任一情况下立即结束：轮到某一方进行回合时，该方无法在落子阶段执行任何合法落子操作；
任一方宣布投降；双方协商一致主动结束当前对局。

游戏结束时，双方比较棋盘上己方棋子的枚数，由于蓝方后手，蓝方获得9枚棋子补偿。
即比较 $\text{红方实际棋子数 与 蓝方实际棋子数}+9$ ，
计分较高的一方获胜。

若游戏因某方投降而结束，则无需进行棋子数量比较，投降方直接判负。

%------------------------------------------------
\section{总结性说明}

炮棋是一种围绕“扩张—成炮—打炮—围捕—反制”展开的双人策略棋类游戏。
玩家首先通过边相邻落子与有限升级扩张己方阵地，再通过形成炮管发动打炮，对棋盘造成前向冲击与内部升级。
打炮既可能强化己方结构，也可能摧毁、削弱甚至翻转敌方棋子，同时还可能形成新的炮管。
除此之外，玩家还可以借助包围和等级优势在吃子阶段实现局部转换占领。
由于打炮与吃子都会改变棋盘状态，而棋盘变化又会不断引发新的炮管与新的可吃子，因此一次落子后，局面常常会发生持续往复的连锁结算。
整局游戏最终在某方无法落子、投降或双方协商结束时终止，并以棋子数量作为胜负依据，蓝方因后手而获得固定九子的结算补偿。
\end{document}